\documentclass{article}
\usepackage[T1]{fontenc}
\usepackage{lmodern}
\usepackage{graphicx}
\usepackage{amsmath,amssymb}
\usepackage[a4paper,margin=2cm]{geometry}
\usepackage[absolute,overlay]{textpos}
\setlength{\TPHorizModule}{1mm}
\setlength{\TPVertModule}{1mm}
\usepackage[indonesian]{babel}
\usepackage{hyperref}
\setlength{\emergencystretch}{2em}
% unit: Halaman 3

\begin{document}

% === Halaman 3 (llm) ===

\section*{Teori Aljabar Abstrak}

\begin{itemize}
    \item Sebelum membahas ECC, perlu dipahami beberapa konsep aljabar abstrak yang mendasarinya.
\end{itemize}

\begin{itemize}
    \item Konsep aljabar abstrak:
    \begin{enumerate}
        \item Grup (\textit{group})
        \item Cincin (\textit{ring})
        \item Medan (\textit{field})
        \item Galois field
    \end{enumerate}
\end{itemize}

\end{document}
